% GENERATED FILE. Edit canonical lessons, then run npm run generate:books.
% canonical-source-hash: fnv1a64:6f84a5ccce445ccb
% canonical-lessons: TE-C106-avakasam, TE-C106-parikaram, TE-C106-asti, TE-C106-yajamani, TE-C106-balam

\chapter{Chance, Tool, Property, Owner, Strength}
\label{ch:te-chance-tool-property-owner-strength}
\hlchaptermodality{106}

\begin{chapteropening}
\textbf{By the end of this chapter:} \emph{I can say a chance, a tool, property, an owner and strength.}

Complete the last lesson of chapter 106: I can say a chance, a tool, property, an owner and strength.
\end{chapteropening}

\section[\texorpdfstring{\te{అవకాశం} (avakāśaṁ)}{అవకాశం (avakāśaṁ)}]{\te{అవకాశం} (avakāśaṁ) — a chance, an opportunity}

\label{lesson:TE-C106-avakasam}



\begin{quote}
Before the new one: say the Telugu for a skill, then the Telugu for a right.
\end{quote}

\subsection*{You'll want to know: \te{అవకాశం}}
\textbf{\te{అవకాశం}} — \emph{avakāśaṁ} — \textquotedblleft{}a chance, an opportunity\textquotedblright{}.

\begin{grammarlens}[title={What you've built}]
One more word for this chapter.
\end{grammarlens}

\subsection*{Your turn}
\begin{itemize}
\raggedright
  \item \emph{Say it:} \emph{avakāśaṁ}
  \item \emph{Say it:} \emph{avakāśaṁ}, once more
  \item \emph{Say it:} say \emph{avakāśaṁ}
  \item \emph{Recall:} say the Telugu for a skill, then the Telugu for a right, then say \emph{avakāśaṁ} again
\end{itemize}

\begin{tcolorbox}[breakable,colback=teal!4,colframe=teal!35!black,title={Before you move on}]
What does \te{అవకాశం} mean? (\textquotedblleft{}A chance, an opportunity\textquotedblright{}.) Say it once more.
\end{tcolorbox}
\section[\texorpdfstring{\te{పరికరం} (parikaraṁ)}{పరికరం (parikaraṁ)}]{\te{పరికరం} (parikaraṁ) — a tool}

\label{lesson:TE-C106-parikaram}



\begin{quote}
Before the new one: say the Telugu for a right, then the Telugu for a chance.
\end{quote}

\subsection*{You'll want to know: \te{పరికరం}}
\textbf{\te{పరికరం}} — \emph{parikaraṁ} — \textquotedblleft{}a tool\textquotedblright{}.

\begin{grammarlens}[title={What you've built}]
One more word for this chapter.
\end{grammarlens}

\subsection*{Your turn}
\begin{itemize}
\raggedright
  \item \emph{Say it:} \emph{parikaraṁ}
  \item \emph{Say it:} \emph{parikaraṁ}, once more
  \item \emph{Say it:} say \emph{parikaraṁ}
  \item \emph{Recall:} say the Telugu for a right, then the Telugu for a chance, then say \emph{parikaraṁ} again
\end{itemize}

\begin{tcolorbox}[breakable,colback=teal!4,colframe=teal!35!black,title={Before you move on}]
What does \te{పరికరం} mean? (\textquotedblleft{}A tool\textquotedblright{}.) Say it once more.
\end{tcolorbox}
\section[\texorpdfstring{\te{ఆస్తి} (āsti)}{ఆస్తి (āsti)}]{\te{ఆస్తి} (āsti) — property}

\label{lesson:TE-C106-asti}



\begin{quote}
Before the new one: say the Telugu for a chance, then the Telugu for a tool.
\end{quote}

\subsection*{You'll want to know: \te{ఆస్తి}}
\textbf{\te{ఆస్తి}} — \emph{āsti} — \textquotedblleft{}property\textquotedblright{}.

\begin{grammarlens}[title={What you've built}]
One more word for this chapter.
\end{grammarlens}

\subsection*{Your turn}
\begin{itemize}
\raggedright
  \item \emph{Say it:} \emph{āsti}
  \item \emph{Say it:} \emph{āsti}, once more
  \item \emph{Say it:} say \emph{āsti}
  \item \emph{Recall:} say the Telugu for a chance, then the Telugu for a tool, then say \emph{āsti} again
\end{itemize}

\begin{tcolorbox}[breakable,colback=teal!4,colframe=teal!35!black,title={Before you move on}]
What does \te{ఆస్తి} mean? (\textquotedblleft{}Property\textquotedblright{}.) Say it once more.
\end{tcolorbox}
\section[\texorpdfstring{\te{యజమాని} (yajamāni)}{యజమాని (yajamāni)}]{\te{యజమాని} (yajamāni) — an owner}

\label{lesson:TE-C106-yajamani}



\begin{quote}
Before the new one: say the Telugu for a tool, then the Telugu for property.
\end{quote}

\subsection*{You'll want to know: \te{యజమాని}}
\textbf{\te{యజమాని}} — \emph{yajamāni} — \textquotedblleft{}an owner\textquotedblright{}.

\begin{grammarlens}[title={What you've built}]
One more word for this chapter.
\end{grammarlens}

\subsection*{Your turn}
\begin{itemize}
\raggedright
  \item \emph{Say it:} \emph{yajamāni}
  \item \emph{Say it:} \emph{yajamāni}, once more
  \item \emph{Say it:} say \emph{yajamāni}
  \item \emph{Recall:} say the Telugu for a tool, then the Telugu for property, then say \emph{yajamāni} again
\end{itemize}

\begin{tcolorbox}[breakable,colback=teal!4,colframe=teal!35!black,title={Before you move on}]
What does \te{యజమాని} mean? (\textquotedblleft{}An owner\textquotedblright{}.) Say it once more.
\end{tcolorbox}
\section[\texorpdfstring{\te{బలం} (balaṁ)}{బలం (balaṁ)}]{\te{బలం} (balaṁ) — strength}

\label{lesson:TE-C106-balam}



\begin{quote}
Before the new one: say the Telugu for property, then the Telugu for an owner.
\end{quote}

\subsection*{You'll want to know: \te{బలం}}
\textbf{\te{బలం}} — \emph{balaṁ} — \textquotedblleft{}strength\textquotedblright{}.

\begin{grammarlens}[title={What you've built}]
One more word for this chapter.
\end{grammarlens}

\subsection*{Your turn}
\begin{itemize}
\raggedright
  \item \emph{Say it:} \emph{balaṁ}
  \item \emph{Say it:} \emph{balaṁ}, once more
  \item \emph{Say it:} say \emph{balaṁ}
  \item \emph{Recall:} say the Telugu for property, then the Telugu for an owner, then say \emph{balaṁ} again
\end{itemize}

\begin{tcolorbox}[breakable,colback=teal!4,colframe=teal!35!black,title={Before you move on}]
What does \te{బలం} mean? (\textquotedblleft{}Strength\textquotedblright{}.) Say it once more.
\end{tcolorbox}
